% =====================================================================
%  Part I -- Astrophysical context.
%  Source: tier0.md I.1 (954-1007), I.2 (1009-1101), VI.6 (3070-3092),
%          I.9 items 1-3 (1427-1432), VI.7 item 6 (3106-3107, relocated).
% =====================================================================
\section{Astrophysical context}
\label{part:context}

Each link in this chain is a physical claim you should be able to defend.

\begin{aside}
\textbf{Notation, taken on credit.} This section is written throughout in \Ye{},
the \emph{electron fraction}: electrons per baryon,
$\Ye = \sum_i Z_i X_i / A_i$, equal to $0.5$ for matter with equal numbers of
protons and neutrons and falling as electrons are captured. It is derived
properly, with its lepton bookkeeping, in \S\ref{sec:ye}. Take it on credit
until then; everything below turns on it.
\end{aside}

% ---------------------------------------------------------------------
\subsection{The last day of a massive star}
\label{sec:lastday}

\subsubsection{The burning sequence}
\label{sec:burningsequence}

A star above ${\sim}10$--11\,\Msun{} burns through the full sequence of fuels
to an iron core, each stage ignited when the ash of the previous one contracts
and heats until the next Coulomb barrier becomes penetrable. (Stars of
${\sim}8$--11\,\Msun{} ignite carbon but generally stop there, ending as ONeMg
cores rather than completing the chain to silicon.) Temperatures, reactions,
and durations below follow \citet{whw2002} --- durations shorten steeply with
mass:

\begin{center}
\footnotesize
\begin{tabularx}{\textwidth}{@{}l l R{1.071} l R{0.929}@{}}
\toprule
\textbf{Stage} & \textbf{Core $T$ (GK)} & \textbf{Principal reactions} &
\textbf{Ash} & \textbf{Duration (15--25\,\Msun{} models)} \\
\midrule
H  & 0.03 & pp chains, CNO cycle & \nuc{He}{4} & ${\sim}10^{7}$\,yr \\
He & 0.2  & $3\alpha \to \nuc{C}{12}$; $\nuc{C}{12}(\alpha,\gamma)\nuc{O}{16}$
          & \nuc{C}{12}, \nuc{O}{16} & ${\sim}10^{6}$\,yr \\
C  & 0.8  & $\nuc{C}{12}(\nuc{C}{12},\alpha)\nuc{Ne}{20}$,
            $(\nuc{C}{12},p)\nuc{Na}{23}$
          & \nuc{Ne}{20}, \nuc{Na}{23}, \nuc{Mg}{24} & ${\sim}10^{3}$\,yr \\
Ne & 1.5  & $\nuc{Ne}{20}(\gamma,\alpha)\nuc{O}{16}$ then
            $\nuc{Ne}{20}(\alpha,\gamma)\nuc{Mg}{24}$
          & \nuc{O}{16}, \nuc{Mg}{24} & ${\sim}1$\,yr \\
O  & 2.0  & $\nuc{O}{16}(\nuc{O}{16},\alpha)\nuc{Si}{28}$,
            $(\nuc{O}{16},p)\nuc{P}{31}$
          & \nuc{Si}{28}, \nuc{S}{32}
          & ${\sim}$months (up to ${\sim}$yr with convection) \\
\textbf{Si} & \textbf{3--4} & \textbf{photodisintegration rearrangement}
          & \textbf{Fe peak} & \textbf{${\sim}1$\,day--2 weeks} \\
\bottomrule
\end{tabularx}
\end{center}

Note the qualitative change at neon: at neon, and again at silicon,
\textbf{photodisintegration initiates the burning} --- oxygen burning, in
between, is still true heavy-ion fusion. The Coulomb barrier for
$\nuc{Si}{28} + \nuc{Si}{28}$ is prohibitive ($Z_1 Z_2 = 196$), so silicon does
not fuse with itself. Instead $\gamma$-rays in the thermal tail knock
$\alpha$-particles, protons, and neutrons off the least-bound nuclei, and those
light particles are captured by others. \textbf{Silicon burning is a
rearrangement, not a fusion reaction.}

This single fact explains most of the project's structure: rearrangement
proceeds through many competing forward/reverse channels near balance, which is
where cancellation ($\kappa$, \S\ref{sec:cancellation}), quasi-equilibrium
clustering (\S\ref{sec:topology}), and stiffness (\S\ref{sec:stiffness}) all
come from.

\subsubsection{Why the acceleration}
\label{sec:acceleration}

From carbon burning onward, energy leaves the core as \textbf{neutrinos}, not
photons \citep{limongi2017,jose2011}. Photons diffuse --- the radiative
diffusion time through a stellar interior is enormous, which is what makes
hydrogen burning last $10^{7}$ years. Neutrinos free-stream at essentially the
speed of light (\S\ref{sec:freestreaming}, below).

The dominant thermal loss is pair annihilation, whose volumetric emissivity
rises roughly as $\boldsymbol{T^{9}}$ here (the relativistic regime,
$\Tn \gtrsim 2$--3; \S\ref{sec:nuchannels}). A core that contracts and heats
loses energy catastrophically faster, so it must burn faster to compensate.
Each successive stage is shorter by a large factor, and the sequence terminates
at the iron peak where no exothermic fusion remains.

\subsubsection{Why this is the regime to emulate}
\label{sec:whyregime}

Silicon burning concentrates the computational pain: highest temperatures
(fastest rates, stiffest equations), largest networks (the iron peak needs many
species), shortest timesteps --- all in the phase that determines what the star
leaves behind, and where a stellar-evolution code spends a disproportionate
fraction of total runtime inside the network.

% ---------------------------------------------------------------------
\subsection{What actually determines whether the star explodes}
\label{sec:explodability}

At the end of Si burning the core is an iron sphere supported by
\textbf{degenerate electron pressure}, which carries ${\sim}93\%$ of the total
at box-centre conditions (\S\ref{sec:boxcentre}).

Take as given, and verify in \S\ref{part:substrate}, that at these
densities the electrons are relativistic and degenerate and carry that
${\sim}93\%$ share; the derivation cannot be brought forward because it is
evaluated at box-centre conditions that are only fixed in
\S\ref{sec:regimebox}.

\subsubsection{Derivation --- \texorpdfstring{$M_{\mathrm{ch}} \propto Y_e^2$}{Mch propto Ye2}}
\label{sec:mchderivation}

For an ultra-relativistic degenerate electron gas, pressure depends only on
electron number density:
%
\begin{equation}
P = \frac{(3\pi^2)^{1/3}\,\hbar c\, n_e^{4/3}}{4}
\label{eq:urel-eos}
\end{equation}
%
Charge neutrality ties $n_e$ to matter density through \Ye{} (derived in
\S\ref{sec:ye}):
%
\begin{equation*}
n_e = \rho N_A Y_e \qquad \implies \qquad P = K\,(\rho Y_e)^{4/3}
\end{equation*}
%
A polytrope with $\gamma = 4/3$, index $n = 3$. The Lane--Emden solution for
$n = 3$ has the property that total mass is \textbf{independent of central
density}:
%
\begin{equation*}
M = 4\pi M_3\left(\frac{K}{\pi G}\right)^{3/2}, \qquad M_3 \approx 2.018
\end{equation*}
%
Since $K \propto Y_e^{4/3}$ and $M \propto K^{3/2}$:
%
\begin{equation}
M_{\mathrm{ch}} \propto Y_e^2 \qquad \text{numerically} \qquad
M_{\mathrm{ch}} \approx 5.83\,Y_e^2\,\Msun
\label{eq:mch}
\end{equation}
%
(Lane--Emden constant $\xi_1^2|\theta'| = 2.01824$; the same derivation appears
in \citealt{suwa2018}, whose $M_{\mathrm{Ch}0} = 1.46\,\Msun\,(\Ye/0.5)^2$ is
this formula.)

\textbf{Check:} $\Ye = 0.5$ gives $1.457\,\Msun$ --- the textbook value.
$\Ye = 0.45$ gives $1.18\,\Msun$, a \textbf{19\% reduction} in supportable mass
from a 10\% change in one composition scalar.

\subsubsection{The thermal correction}
\label{sec:thermalcorrection}

\begin{equation*}
M_{\mathrm{ch,eff}} \approx
5.83\,Y_e^2\left[1 + \left(\frac{s_e}{\pi Y_e}\right)^{\!2}\right]\Msun
\end{equation*}

where $s_e$ is the electron entropy per baryon in units of $k$. This is the
first-order form of \citet{baron1990}, whose exact bracket is
$[1 + (2/3)(s_e/\pi Y_e)^2]^{3/2}$ (see also \citealt{timmes1996};
\citealt{suwa2018}). So \textbf{entropy also raises the effective
Chandrasekhar mass} --- a hotter core supports more. \Ye{} remains the dominant
lever and the one nuclear burning directly controls, but the picture is
two-parameter.

\subsubsection{The collapse trigger}
\label{sec:collapsetrigger}

Collapse begins when pressure support fails, via two coupled runaways:

\begin{enumerate}
\item \textbf{Photodisintegration of the iron peak.} Once the core reaches
      $\boldsymbol{T \gtrsim 7\text{--}10}$~\textbf{GK} (\citealt{janka2012},
      after B$^2$FH; \citealt{b2fh1957}) --- above the regime box, and strongly
      density-dependent --- $\nuc{Fe}{56}(\gamma,\alpha)$-type reactions
      dismantle the Fe peak back to $\alpha$-particles and free nucleons. This
      is strongly \textbf{endothermic}: $\nuc{Fe}{56} \to 13\alpha + 4n$ costs
      124.4~MeV (AME2020 mass excesses; the convention is
      \S\ref{sec:massexcess}), i.e.\ 2.22~MeV/nucleon, consuming thermal energy
      and reducing pressure.

      \begin{warn}
      \textbf{Do not put this turn-on at ${\sim}5$~GK.} Si burning
      \emph{produces} the Fe peak at 3--4~GK (\S\ref{sec:lastday}), and the
      box's own upper edge of 7.9~GK is characterized as full NSE that is
      still Fe-peak dominated (\S\ref{sec:regimebox}). A dismantling threshold
      at 5~GK would contradict both. The Fe peak survives NSE across the whole
      box; it is collapse --- not silicon burning --- that reaches the
      temperatures where it does not.
      \end{warn}

\item \textbf{Electron capture.} Free protons liberated by (1), and Fe-peak
      nuclei directly, capture electrons. Each capture removes a
      pressure-supporting electron \emph{and} emits a neutrino that leaves.
      \Ye{} drops, so $M_{\mathrm{ch}}$ drops, so the marginally stable core no
      longer is.
\end{enumerate}

These reinforce each other, and collapse becomes dynamical on a free-fall
timescale of ${\sim}0.1$~s ($t_{\mathrm{coll}} \approx 0.21/\sqrt{\bar\rho_8}$~s
--- \citealt{janka2012}).

\subsubsection{The observable end of the chain}
\label{sec:observableend}

\Ye{} at collapse onset sets: the iron core mass, the bounce shock location,
how much material the shock must traverse, and therefore whether the explosion
succeeds. Downstream it sets ejecta neutron-richness and hence observed
abundance patterns. The \textbf{compactness parameter}
$\xi_M = (M/\Msun)/(R(M)/1000~\mathrm{km})$ at $M = 2.5\,\Msun$, evaluated at
core bounce \citep{oconnor2011} --- the standard explodability predictor ---
is set by exactly this core structure (\S\ref{sec:obsgrounding}, item~3, for
how it is measured; the definition is deliberately restated there, where the
measurement chain rather than the physics is the subject).

\textbf{This is why every accuracy gate in this repository is written in
\Ye{}}, and why the invariants say the weak sector may never be masked, frozen,
or projected away. An emulator that conserves mass and charge perfectly but
gets \Ye{} wrong has failed at the only thing that mattered.

% ---------------------------------------------------------------------
\subsection{The observational grounding --- how we know \texorpdfstring{\Ye{}}{Ye} matters}
\label{sec:obsgrounding}

Chains of inference worth being able to state:

\begin{enumerate}
\item \textbf{\nuc{Ni}{56} masses from light curves.} The decay chain
      $\nuc{Ni}{56} \to \nuc{Co}{56} \to \nuc{Fe}{56}$ (half-lives 6.08~d,
      77.2~d --- ENSDF) powers the tail of a supernova light curve
      \citep{arnett1982,diehl2015}. Fitting the tail gives the ejected
      \nuc{Ni}{56} mass directly. That mass depends on how much material burned
      to NSE (\S\ref{sec:chempotential}) and at what \Ye{}.

\item \textbf{Abundance patterns in metal-poor stars.} Fe-peak element ratios
      (Ni/Fe, Cr/Fe, Mn/Fe) are sensitive to \Ye{} in the explosive burning
      region. Observed ratios constrain the \Ye{} of the ejected burning region
      (e.g.\ \citealt{jerkstrand2015} for Ni/Fe).

\item \textbf{Explodability studies.} The compactness parameter
      $\xi_M = (M/\Msun)/(R(M)/1000~\mathrm{km})$ at $M = 2.5\,\Msun$,
      evaluated at bounce, correlates with whether simulations explode
      \citep{oconnor2011}. It is set by the core structure at collapse --- set
      in turn by the Si-burning history.

\item \textbf{SN 1987A.} Direct neutrino detection confirmed the core-collapse
      energetics picture and the neutrino emission scale
      \citep{hirata1987,bionta1987}.
\end{enumerate}

This is what makes the accuracy targets \emph{meaningful} rather than
self-referential: there is a measurement at the end of the chain.

% ---------------------------------------------------------------------
\subsection*{Self-check --- astrophysical context}
\addcontentsline{toc}{subsection}{Self-check --- astrophysical context}
\label{sec:selfcheck-context}

\begin{selfcheck}
\item Explain why the burning-stage durations shorten by orders of magnitude,
      naming the mechanism and the temperature scaling.
\item Derive $M_{\mathrm{ch}} \propto Y_e^2$ from the ultra-relativistic
      degenerate EOS and the $n = 3$ polytrope. Where does the
      density-independence of $M$ come from?
\item State the two coupled runaways that trigger collapse and explain how each
      reduces pressure support.
\item Trace the inference chain from ``an emulator gets \Ye{} wrong by 1\%'' to
      ``an observable changes'', naming each physical link.
      \emph{(Relocated here from the source document's \S VI.7 item~6, which
      tested \S VI.6 --- the material now in \S\ref{sec:obsgrounding}.)}
\end{selfcheck}
